\documentclass{article}
\usepackage{fancyhdr}
\usepackage{ctex}
\usepackage{listings}
\usepackage[a4paper, body={18cm,22cm}]{geometry}
\usepackage{amsmath,amssymb,amstext,wasysym,enumerate,graphicx}
\usepackage{float,abstract,booktabs,indentfirst,amsmath}
\usepackage{multirow}
\usepackage{enumitem}
\usepackage{listings}
\usepackage{xcolor}
\usepackage{tabularx}
\usepackage[most, skins, breakable, theorems]{tcolorbox}
\usepackage{accsupp}
\usepackage[bottom]{footmisc}
\usepackage{subcaption}
\usepackage{logicproof}
% \usepackage[backend=biber,style=numeric]{biblatex}
\usepackage[xetex]{hyperref}
\usepackage{fontspec}
\usepackage{amsthm}
\newtheorem{theorem}{Theorem}
\usetikzlibrary{arrows.meta}
\newcommand\emptyaccsupp[1]{\BeginAccSupp{ActualText={}}#1\EndAccSupp{}}
\setlength{\parindent}{2em}
\setmonofont{Consolas}
\setCJKmonofont{黑体}
% \setmainfont{Times New Roman}
\hypersetup{CJKbookmarks=true,colorlinks=true,citecolor=blue,%
            linkcolor=blue,urlcolor=blue,bookmarksnumbered=true,%
            bookmarksopen=true,breaklinks=true}
\lstset{
    % language = C,
    xleftmargin = 3em,xrightmargin = 3em, aboveskip = 1em,
	backgroundcolor = \color{white}, % 背景色
	basicstyle = \small\ttfamily, % 基本样式 + 小号字体
	rulesepcolor= \color{gray}, % 代码块边框颜色
	breaklines = true, % 代码过长则换行
	numbers = left, % 行号在左侧显示
	numberstyle = \small\emptyaccsupp, % 行号字体
    numbersep = -14pt, 
    keywordstyle=\color{purple}\bfseries, % 关键字颜色
    commentstyle =\color{red!50!green!50!blue!60}, % 注释颜色
    stringstyle = \color{red}, % 字符串颜色
    morekeywords={ASSERT, int64_t, uint32_t},
	frame = single, 
	showspaces = false, % 不显示空格
    showstringspaces = false,
	columns = fixed, % 字间距固定
    literate=
        {^-}{{{\color{black}{\color{red}-}}\colorbox{red!30}{\phantom{XX}}}}{1}
        {^+}{{{\color{black}{\color{green}+}}\colorbox{green!30}{\phantom{XX}}}}{1},
}

\raggedbottom

\title{\heiti{计算机逻辑基础}}
\author{第七次作业 \\ 
武泽恺 10225101429
}
\date{2025年12月27日}

\begin{document}
\maketitle

Verify the following Hoare triple is partial correct:

\begin{itemize}
    \item \textbf{Precondition:} $\{n > 0 \land \forall i. 0 \le i < n \to a[i] \in \mathbb{N}\}$
    \item \textbf{Program:} Sort
    \item \textbf{Postcondition:} $\{\forall i, j. 0 \le i \le j < n \to a[i] \le a[j]\}$
\end{itemize}

\begin{lstlisting}[language=c]
    i = 0;
    while (i < n-1) {
        j = 0;
        while (j < n-i-1) {
            if (a[j] > a[j+1]) {
                temp = a[j];
                a[j] = a[j+1];
                a[j+1] = temp;
            }
            j = j+1;
        }
        i = i+1;
    }
\end{lstlisting}

我们定义以下循环不变式：
\[
\begin{aligned}
\text{inv1} & =_{df} \forall x, y. x \in [n - i - 1 .. n - 2] \land y \in [x + 1 .. n - 1] \to a[x] \le a[y] \\
\text{inv2} & =_{df} a[j] = \max(a[0 .. j])
\end{aligned}
\]
也就是说，在外层循环不变式inv1下，数组$a[n-i-1 .. n-1]$是有序的；
在内层循环不变式inv2下，$a[0 .. j]$中的最大值存储在$a[j]$中。

\begin{proof}
\noindent\begin{tabular}{lll}
\hline
{Line} & {Assertion} & {Rule} \\ \hline
& $\{n > 0 \land \forall i. 0 \le i < n \to a[i] \in \mathbb{N}\}$ & pre \\ \hline
& $\{\forall x, y. x \in [n - 1 .. n - 2] \land y \in [x + 1 .. n - 1] \to a[x] \le a[y]\}$ & implied \\ \hline
1 & $i = 0;$ & \\ \hline
& $\{inv1\}$ & assignment \\ \hline
2 & {while} $(i < n - 1)$ \{ & \\ \hline
& $\{inv1 \land i < n - 1\}$ & while \\ \hline
& $\{inv1 \land i < n - 1 \land a[0] = \max(a[0 .. 0])\}$ & implied \\ \hline
3 & $j = 0;$ & \\ \hline
& $\{inv1 \land i < n - 1 \land a[j] = \max(a[0 .. j])\}$ & assignment \\ \hline
& $\{inv1 \land i < n - 1 \land inv2\}$ & implied \\ \hline
4 & {while} $(j < n - i - 1)$ \{ & \\ \hline
& $\{inv1 \land i < n - 1 \land inv2 \land j < n - i - 1\}$ & while \\ \hline
5 & {if} $(a[j] > a[j + 1])$ \{ & \\ \hline
& $\{inv1 \land i < n - 1 \land inv2 \land a[j] > a[j + 1]\}$ & if \\ \hline
6 & $temp = a[j];$ & \\ \hline
& $\{inv1 \land i < n - 1 \land inv2 \land temp > a[j + 1]\}$ & assignment \\ \hline
7 & $a[j] = a[j + 1];$ & \\ \hline
& $\{inv1 \land i < n - 1 \land inv2 \land temp > a[j]\}$ & assignment \\ \hline
8 & $a[j + 1] = temp;$ & \\ \hline
& $\{inv1 \land i < n - 1 \land a[j] < a[j + 1]\}$ & assignment \\ \hline
9 & \} & \\ \hline
& $\{inv1 \land i < n - 1 \land a[j + 1] = \max(a[0 .. j + 1]) \land a[j] \le a[j + 1]\}$ & if \\ \hline
10 & $j = j + 1;$ & \\ \hline
& $\{inv1 \land i < n - 1 \land a[j] = \max(a[0 .. j])\}$ & assignment \\ \hline
11 & \} & \\ \hline
& $\{inv1 \land i < n - 1 \land a[j] = \max(a[0 .. j]) \land j \ge n - i - 1\}$ & while \\ \hline
& $\{\forall x, y. x \in [n - i - 1 .. n - 2] \land y \in [x + 1 .. n - 1] \to a[x] \le a[y] \dots\}$ & implied \\ \hline
& $\{\forall x, y. x \in [n - i - 2 .. n - 2] \land y \in [x + 1 .. n - 1] \to a[x] \le a[y] \land i < n - 1\}$ & implied \\ \hline
12 & $i = i + 1;$ & \\ \hline
& $\{\forall x, y. x \in [n - i - 1 .. n - 2] \land y \in [x + 1 .. n - 1] \to a[x] \le a[y] \land i < n\}$ & assignment \\ \hline
& $\{\forall x, y. x \in [n - i - 1 .. n - 2] \land y \in [x + 1 .. n - 1] \to a[x] \le a[y]\}$ & implied \\ \hline
13 & \} & \\ \hline
& $\{inv1 \land i \ge n - 1\}$ & while \\ \hline
& $\{\forall x, y. x \in [0 .. n - 2] \land y \in [x + 1 .. n - 1] \to a[x] \le a[y]\}$ & implied \\ \hline
& $\{\forall i, j. 0 \le i \le j < n \to a[i] \le a[j]\}$ & post \\ \hline
\end{tabular}
\end{proof}


\end{document}